\chapter{Appendix F: Conventions and Normalizations}
\label{app:conventions}

\section{Metric and Geometric Normalizations}
To ensure consistency across coupling regimes and avoid the "floating constant" ambiguity noted by referees, we fix the following conventions for the Lie group $G = SU(N)$.

\subsection{Lie Algebra Metric}
We equip the Lie algebra $\mathfrak{su}(N)$ with the Killing form scaled to match the trace in the fundamental definition:
\begin{equation}
    \langle X, Y \rangle = - \frac{1}{2} \text{Tr}(XY), \quad X, Y \in \mathfrak{su}(N)
\end{equation}
This normalization ensures that for $SU(2)$, the Pauli matrices $\sigma_i$ satisfy $\langle i\sigma_a, i\sigma_b \rangle = \delta_{ab}$.

\subsection{Haar Measure Probability}
The Haar measure $dU$ on $SU(N)$ is always normalized to be a probability measure:
\begin{equation}
    \int_{SU(N)} dU = 1
\end{equation}

\subsection{Ricci Curvature}
The Ricci curvature tensor is defined via the Bochner formula for the laplacian $\Delta = \sum_i X_i^2$:
\begin{equation}
    \frac{1}{2} \Delta |\nabla f|^2 = |\nabla^2 f|^2 + \langle \nabla f, \nabla \Delta f \rangle + \text{Ric}(\nabla f, \nabla f)
\end{equation}
With our metric normalization, the Ricci curvature of $SU(N)$ is strictly positive:
\begin{equation}
    \text{Ric}_{SU(N)} = \frac{N}{2} \cdot g_{\mathfrak{su}(N)}
\end{equation}
With this normalization of the metric/Laplacian, the Bakry--\'Emery criterion yields the Haar Log-Sobolev constant
\begin{equation}
    \rho_{Haar} = \frac{N}{2}.
\end{equation}
\emph{Translation note.} Some references use a different normalization of the Laplacian (or a different choice of inner product on $\mathfrak{su}(N)$), in which case the numerical value of $\rho_{Haar}$ changes by a \emph{universal constant factor}. Throughout this manuscript, we use the convention fixed here and in Appendix~B, so $\rho_{Haar}$ always scales linearly in $N$.

\section{Lattice Hamiltonian Normalization}
The lattice Hamiltonian $H$ is defined from the transfer matrix $\hat{T}$ in the temporal direction.
\begin{equation}
    \hat{T} = e^{-a H}, \quad H \ge 0
\end{equation}
The Hamiltonian density is intensive. When comparing to the Glauber dynamics generator $\mathcal{L}$, we use the normalized Dirichlet form:
\begin{equation}
    \mathcal{E}(f,g) = \lim_{\Lambda \to \infty} \frac{1}{|\Lambda|} \sum_{x \in \Lambda} \langle \nabla_x f, \nabla_x g \rangle_{L^2}
\end{equation}
This ensures that spectral gaps $\gamma$ do not scale with volume $V$, avoiding the $1/V$ vanishing gap paradox.
